% ============================================================
% QUANTUM SECTOR — ECT
% File: qs.tex
% Version: 4  (2026-03-17, §7.3 approved, §7.4 added)
%
% VERSION HISTORY:
%   v1:  §7.1 initial draft.
%   v2:  §7.1 polished (GPT round 1); §7.2 added.
%   v3:  §7.2 Fix A: u_0^2 -> K_theta^{eff}. Fix B: fake cite removed.
%        §7.3 Pre-quantum action scale and the status of hbar.
%   v4:  §7.3 approved by GPT. §7.4 Wave dynamics, Wick rotation,
%        and Schrodinger-type evolution added.
%          - Phase action recalled via \eqref (no duplicate equation)
%          - \subsubsection* used for major structural blocks in §7.4
%            (longer subsection, appropriate); \paragraph for sub-steps
%          - eq:qs_wick labelled (Wick rotation w->-ic_*t, no existing
%            label in ECT_preprint; noted as established in ECT basics)
%          - eq:qs_wave_ansatz WKB-like structure noted explicitly
%          - eq:qs_schrodinger_type: S_0 appears in role of hbar;
%            clear Level B status
%          - Logic chain diagram labelled eq:qs_logic_chain
%          - Bridge leads to sec:qs_ccr (canonical structure)
%
% DISCIPLINE:
% (1) Every claim: Level A / B / EFT / matching / open.
% (2) Two parallel lines: mathematical + physical.
% (3) No circular use of hbar, Schrodinger, Born rule before derivation.
% (4) Deferred basics DEVELOPED here:
%     sec:coh_branch, sec:hbar_status, sec:arrow_of_time,
%     app:ohmic_markov, app:Sloop_calc, sec:defect_sectors.
% (5) No unjustified identifications.
% (6) Every strong claim \ref{} or \cite{}.
% (7) Explicit Level A/B/C for every non-trivial statement.
%
% INTEGRATION NOTE:
%   Replaces: legacy sec:emergent_principles + sec:quantum_vacuum.
%   Main label: sec:quantum (use sec:quantum_new during development).
%   Subsection labels: sec:qs_intro, sec:qs_phase_dynamics,
%     sec:qs_hbar, sec:qs_wave, sec:qs_ccr, sec:qs_decoherence,
%     sec:qs_born, sec:qs_dark_sector, sec:qs_defects, sec:qs_status.
% ============================================================

\section{Quantum Sector}
\label{sec:quantum_new}
% NOTE: upon integration becomes \label{sec:quantum}

% ============================================================
\subsection{Quantum Sector as the development of the
coherent-topological branch}
\label{sec:qs_intro}
% ============================================================

\textit{Connection to ECT basics:
Sections~\ref{sec:branching}, \ref{sec:coh_branch},
\ref{sec:hbar_status}.
Status: conceptual bridge; all structural ingredients were
prepared in ECT basics and are now developed dynamically.}

\medskip
After ECT basics and Macroscopic Physics, the logical structure
of ECT has separated into two effective branches from the same
ordered condensate.
The first branch---developed in Macroscopic Physics
(Section~\ref{sec:macroscopic})---describes those observables
governed by the phase-insensitive, coarse-grained geometric
response: long-wavelength metric fluctuations, gravitational
stiffness, cosmological dynamics, and galactic phenomenology.
In that branch the global phase of the effective order parameter
is irrelevant to the observables of interest and can be averaged away.

The second branch---developed here---treats the same ordered
condensate through observables that \emph{remain sensitive} to
global phase structure, winding sectors, loop configurations,
and coherent correlators.
This is the coherent-topological branch of ECT, and its systematic
development is the task of the Quantum Sector.

\paragraph{Why a separate Quantum Sector is needed.}
The geometric branch is sufficient whenever local phase information
is washed out under coarse-graining.
It is not sufficient for phenomena in which the global phase of
the effective order parameter, the topology of loop sectors, and the
coherent structure of amplitudes remain physically active.
Those phenomena require a different effective description---one that
retains global phase data rather than discarding it.
The Quantum Sector provides this description.

\paragraph{Primary variables of the coherent-topological branch.}
In the coherent branch the ordered condensate is described by
variables that retain global phase information.
The primary objects are:
\begin{itemize}
  \item the effective order parameter
    \begin{equation}
      \Phi_{\rm eff} = \rho\,e^{i\theta}
      \label{eq:qs_order_parameter}
    \end{equation}
  \item the phase variable $\theta$, which is \emph{not}
    coarse-grained away;
  \item the winding quantisation condition,
    \begin{equation}
      \oint_\gamma \partial_A\theta\,dX^A = 2\pi n,
      \qquad n\in\mathbb{Z}
      \label{eq:qs_winding}
    \end{equation}
  \item the pre-quantum action scale $S_0$, introduced through
    the minimal nontrivial loop action
    \begin{equation}
      S_{\rm loop,min} = 2\pi S_0
      \label{eq:qs_Sloop}
    \end{equation}
  \item coherent correlators and loop amplitudes of the ordered sector.
\end{itemize}
These objects were already structurally prepared in ECT basics:
the coherent-topological branch was identified
(Section~\ref{sec:coh_branch}), the pre-quantum scale $S_0$ was
introduced (Section~\ref{sec:hbar_status}), and the existence
of winding sectors was established~(P5).
None of them was fully developed dynamically; that development
belongs to the present section.

\paragraph{What ECT basics already established.}
The following structural ingredients are available before
entering the Quantum Sector:
\begin{enumerate}[label=(\roman*)]
  \item The ordered phase admits a genuine phase-sensitive sector,
    distinct from the geometric branch~(Level~A within the
    broken-phase EFT).
  \item Topologically quantised winding sectors exist in the
    coherent order parameter as a consequence of P5~(Level~A).
  \item A pre-quantum action scale $S_0$ is structurally
    introduced by the loop quantisation condition~\eqref{eq:qs_Sloop}
    (Level~A for the structural existence of the scale;
    its explicit first-principles evaluation is deferred,
    see Appendix~\ref{app:Sloop_calc}).
  \item The variational/Noether foundation is common to both
    branches, so the Quantum Sector is not an independent add-on
    but a second effective development of the same condensate
    theory (Level~A: follows from the shared P1--P3 action).
\end{enumerate}

\paragraph{What has not yet been established at this point.}
ECT basics and Macroscopic Physics did \emph{not} yet prove:
\begin{enumerate}[label=(\roman*)]
  \item That $S_0 = \hbar$.
    This is one of the first central tasks of the Quantum Sector:
    it requires an explicit calculation and identification,
    not merely the structural existence of $S_0$.
  \item That canonical commutation relations follow in their
    standard form.
  \item That the Born rule has been derived.
  \item That decoherence, probability interpretation, and
    measurement theory are complete.
  \item That topological defect sectors (vortices, domain walls,
    instantons) have been fully developed.
\end{enumerate}
All these belong to the present section or to its explicitly
listed open problems.

\paragraph{Why the phase is not coarse-grained away here.}
In Macroscopic Physics, local phase information was irrelevant
because the observables under study were governed by the coarse-grained
metric response and long-wavelength stiffness.
In the coherent-topological branch the relevant observables are
different: the global consistency of the phase, the existence of
nontrivial loop sectors, and the quantisation of winding numbers
carry direct physical content.
Therefore the phase variable $\theta$ is not discarded under
coarse-graining; it becomes one of the principal dynamical variables
of the effective description.

\paragraph{Physical meaning of the Quantum Sector.}
The Quantum Sector should \emph{not} be understood as importing
ordinary quantum mechanics into ECT from outside.
Its role is precisely the opposite: to show how phase-sensitive and
topological features of the ordered condensate can \emph{reproduce,
motivate, or constrain} the structures that ordinary quantum theory
usually postulates.
Both branches---geometric and coherent-topological---are consequences
of the same five postulates P1--P5; they differ only in which class
of observables they target.

\paragraph{Logical structure of the Quantum Sector.}
The present section proceeds in the following order:
\begin{enumerate}[label=(\arabic*)]
  \item \textbf{Coherent phase dynamics and loop sectors}
    (Section~\ref{sec:qs_phase_dynamics}):
    the dynamics of $\theta$, loop configurations, and why
    phase dynamics is the correct effective language of the
    coherent branch.
  \item \textbf{Pre-quantum action scale and the status of $\hbar$}
    (Section~\ref{sec:qs_hbar}):
    development of Section~\ref{sec:hbar_status}
    and Appendix~\ref{app:Sloop_calc};
    conditions for $S_0 = \hbar$.
  \item \textbf{Wave dynamics, Wick rotation, and
    Schr\"odinger-type evolution}
    (Section~\ref{sec:qs_wave}):
    coherent amplitude language, Lorentzian continuation
    $w\to -ic_*t$, and the emergence of first-order
    coherent wave evolution.
  \item \textbf{Canonical structure and the status of CCR}
    (Section~\ref{sec:qs_ccr}):
    where and in what sense canonical commutation relations
    appear; what is already established and what is not.
  \item \textbf{Influence functional, decoherence, and the
    arrow of time}
    (Section~\ref{sec:qs_decoherence}):
    development of Section~\ref{sec:arrow_of_time} and
    Appendix~\ref{app:ohmic_markov};
    causal kernel, entropy growth, Ohmic/Markov assumptions,
    classical/quantum split.
  \item \textbf{Probability, Born-type interpretation, and
    measurement status}
    (Section~\ref{sec:qs_born}):
    what can be established; what remains Level~B or open.
  \item \textbf{Dark sector from the coherent-topological branch}
    (Section~\ref{sec:qs_dark_sector}):
    full development of the MP bridge to collective and
    topological dark-sector candidates.
  \item \textbf{Defect sectors beyond the smooth coherent branch}
    (Section~\ref{sec:qs_defects}):
    development of Section~\ref{sec:defect_sectors}:
    vortices, domain walls, instantons; why P5 chooses the smooth
    coherent sector but does not exclude defect ontology.
  \item \textbf{Status table and open problems}
    (Section~\ref{sec:qs_status}).
\end{enumerate}

\paragraph{Bridge from Macroscopic Physics.}
Macroscopic Physics answered the question:
\emph{how does the ordered condensate appear to a large-scale
observer insensitive to coherent phase data?}
The Quantum Sector asks the complementary question:
\emph{what effective dynamics remains when that phase data is
retained rather than averaged away?}
These are not two different theories but two branches of the
same ordered condensate, selected by two different classes of
observables.
The bridge between them is the branching structure established
in ECT basics (Section~\ref{sec:branching}).

\paragraph{Summary.}
The Quantum Sector is the systematic development of the
coherent-topological branch of ECT.
Its task is to turn the structural ingredients prepared in
ECT basics---phase variable $\theta$, winding sectors~(P5),
loop quantisation, and the pre-quantum action scale $S_0$---into
a controlled framework for wave dynamics, decoherence,
probability, and topological dark-sector physics.
All results carry explicit Level~A/B/C markers.

% ============================================================
\subsection{Coherent phase dynamics and loop sectors}
\label{sec:qs_phase_dynamics}
% ============================================================

\textit{Status: Level~A for the existence of the phase variable,
winding sectors, and the topological loop classification inside the
coherent branch; Level~B for the specific effective dynamical closure
used to connect these structures to wave dynamics.
Connection to ECT basics: P5, Section~\ref{sec:hbar_status},
Appendix~\ref{app:Sloop_calc}.}

\medskip
The first task of the Quantum Sector is to identify the correct
effective variables of the coherent-topological branch and to show
why their natural language is the language of phase dynamics.
This is not yet ordinary quantum mechanics.
Rather, it is the phase-sensitive effective description of the ordered
condensate in the regime where coherent loop sectors remain
physically relevant.

\paragraph{Phase as the natural variable of the coherent branch.}
Once the ordered condensate is written as
\begin{equation}
  \Phi_{\rm eff}(X) = \rho(X)\,e^{i\theta(X)},
  \label{eq:qs_polar}
\end{equation}
the phase field $\theta$ becomes physically meaningful whenever
configurations with nontrivial winding are retained.
In the geometric branch such information is averaged away.
In the coherent branch it is exactly this global phase structure that
distinguishes one sector from another.
The phase is not merely a convenient reparametrisation: it is the
variable that remembers the global consistency conditions of the
condensate---single-valuedness, winding, and loop topology.

\paragraph{Reduced phase action.}
At wavelengths large compared to the microscopic core scale, and within
the smooth coherent sector selected by P5, the radial mode $\rho$ is
approximately frozen and the effective dynamics is dominated by the
phase field.
The quadratic phase action in this regime takes the generic form
(see Section~\ref{sec:hbar_status} and Appendix~\ref{app:Sloop_calc}):
\begin{equation}
  \label{eq:qs_phase_action}
  S_\theta
  =
  \frac{K_\theta^{\rm eff}}{2}
  \int d^4X\; g_{\rm eff}^{AB}(X)\,\partial_A\theta\,\partial_B\theta
  + S_{\rm higher},
\end{equation}
where $g_{\rm eff}^{AB}$ is the effective Lorentzian metric tensor of
the broken phase---the same tensor that enters the geometric branch at
macroscopic scales (Section~\ref{sec:lorentz_order_field})---and
$S_{\rm higher}$ denotes higher-gradient and interaction corrections.

Here $K_\theta^{\rm eff}$ denotes the \emph{effective phase stiffness}
of the coherent branch.
At the level of the present closure it is expected to be determined by
the ordered condensate background and ultimately related to the same
broken-phase data that define the ordered vacuum~(the stiffness
parameter $u_0$, the radial mass $m_\sigma$, and the anisotropy
$\alpha$).
However, its explicit first-principles reduction to these primary
parameters is not assumed here: it requires an explicit EFT matching
step in which the heavy radial mode and the orientation fluctuations are
integrated out consistently.
That matching belongs to the coherent-sector EFT program and is
left as an open problem pending the full loop-sector
calculation~(Appendix~\ref{app:Sloop_calc}).
Equation~\eqref{eq:qs_phase_action} should be understood as a
Level~B effective closure, not yet the full microscopic derivation
from bare~P3.

\paragraph{Why this action is natural.}
Equation~\eqref{eq:qs_phase_action} is the phase-sector analogue of
the sigma-model and tensor EFT structures already used in ECT basics
and Macroscopic Physics.
Its role is to capture the lowest-order dynamics of smooth coherent
configurations once the heavy radial structure has been integrated out
and the phase is retained as a genuine dynamical degree of freedom.
The phase action is not an extra postulate but the natural effective
continuation of the broken ordered phase into its coherent-topological
branch.

\paragraph{Winding sectors and topological classification.}
The coherent branch is partitioned into sectors classified by winding
number (cf.~\eqref{eq:qs_winding}):
\begin{equation}
  \label{eq:qs_winding_sector}
  n[\gamma] \;=\; \frac{1}{2\pi}
  \oint_\gamma \partial_A\theta\,dX^A \;\in\; \mathbb{Z}.
\end{equation}
This is a Level~A structural statement: it follows from the
single-valuedness of $\Phi_{\rm eff}$ in the smooth coherent sector
selected by P5.
The trivial sector $n=0$ contains contractible smooth configurations;
nontrivial sectors $n\neq0$ correspond to configurations with
topologically nontrivial phase winding that cannot be eliminated by
a smooth deformation.

\paragraph{Minimal loop action and the appearance of $S_0$.}
The existence of nontrivial winding sectors immediately implies a
minimal nontrivial loop contribution to the action.
The $n=1$ minimal loop carries action
\begin{equation}
  \label{eq:qs_loop_min}
  S_{\rm loop,min} = 2\pi S_0,
\end{equation}
where the pre-quantum action scale $S_0$ depends on the coherent-sector
EFT parameters~(Appendix~\ref{app:Sloop_calc}).
At this stage $S_0$ is \emph{not} yet identified with $\hbar$: what
is established is only that the coherent branch carries a distinguished
action scale~(Level~A for existence; Level~B for explicit value pending
the loop-sector calculation).

\paragraph{Physical meaning of the loop sectors.}
A nontrivial loop sector represents more than a winding number.
Physically, it means that the condensate admits coherent histories that
cannot be continuously deformed into the trivial sector without leaving
the smooth coherent branch.
That is why loop sectors are the natural place where quantum-like
notions of amplitude, interference between sectors, and action
quantisation can first arise in ECT---not from operator axioms, but
from the topology of the ordered condensate itself.

\paragraph{Smooth coherent branch versus defect sectors.}
The present discussion is restricted to the smooth coherent sector
chosen by P5.
Configurations with singular cores, vanishing $\rho$, or
discontinuous orientation data belong to the wider defect sector
and are treated in Section~\ref{sec:qs_defects}.
The role of P5 is not to deny the existence of defects, but to define
the branch in which smooth phase dynamics can first be developed
cleanly.

\paragraph{From loop sectors to wave dynamics.}
Once the condensate admits coherent sectors weighted by their phase
action, the next question is how these sectors combine into an
effective wave description.
This step requires: (a) showing how the phase-sector path integral
leads to a Schr\"odinger-like amplitude language; (b) verifying that
the Wick rotation $w\to -ic_*t$ (established for the geometric branch;
Section~\ref{sec:lorentz_order_field}) carries the phase action into a
standard kinetic term; (c) establishing when $S_0$ becomes
identifiable with $\hbar$.
Those questions belong to Sections~\ref{sec:qs_hbar}
and~\ref{sec:qs_wave}.

\paragraph{Status summary.}
\textit{Established}~(Level~A unless noted):
(i)~the coherent branch is naturally described by $\theta$;
(ii)~smooth sectors are classified by integer winding~(P5);
(iii)~nontrivial sectors imply a distinguished action scale $S_0$.
\textit{Not yet established}:
(iv)~the explicit value of $K_\theta^{\rm eff}$ and $S_0$
~(open, Appendix~\ref{app:Sloop_calc});
(v)~the identification $S_0=\hbar$~(Section~\ref{sec:qs_hbar});
(vi)~the full wave-dynamical superposition rule between sectors;
(vii)~the role of defect sectors~(Section~\ref{sec:qs_defects}).

\paragraph{Summary.}
The coherent-topological branch is naturally organised by phase
dynamics.
Its primary data are the phase field $\theta$, the winding
condition~\eqref{eq:qs_winding_sector}, and the loop-sector
classification of the ordered condensate.
This gives ECT a distinguished coherent action scale and a topological
organisation of physically admissible sectors.
Whether this pre-quantum scale can be identified with the observed
$\hbar$ is addressed next.

% ============================================================
\subsection{Pre-quantum action scale and the status of
  \texorpdfstring{$\hbar$}{hbar}}
\label{sec:qs_hbar}
% ============================================================

\textit{Status: Level~A for the structural existence of a distinguished
pre-quantum action scale $S_0$ in the coherent-topological branch;
Level~B for any identification of that scale with the observed $\hbar$.
Connection to ECT basics: Section~\ref{sec:hbar_status} and
Appendix~\ref{app:Sloop_calc}.}

\medskip
The previous subsection established that the coherent branch of ECT
is organised by nontrivial loop sectors of the phase field and that
these sectors carry a minimal nontrivial action
$S_{\rm loop,min}=2\pi S_0$~\eqref{eq:qs_Sloop}.
The present subsection clarifies the precise status of $S_0$.
The key point is that $S_0$ is already structurally present in ECT,
but its identification with the observed quantum of action $\hbar$ is
\emph{not} automatic.

\paragraph{Why $S_0$ exists at Level~A.}
The existence of $S_0$ follows from three structural facts:
\begin{enumerate}[label=(\roman*)]
  \item the coherent branch retains a genuine phase variable $\theta$;
  \item smooth coherent configurations are partitioned into integer
    winding sectors~(P5);
  \item the first nontrivial loop sector defines a minimal nonzero
    contribution to the action.
\end{enumerate}
Together these imply that the coherent branch possesses a distinguished
action scale associated with topologically nontrivial phase transport.
This is a structural consequence of the coherent-topological sector,
not an imported quantum axiom.

\paragraph{What $S_0$ means physically.}
The quantity $S_0$ should be understood as the intrinsic unit of action
prepared by the topology of the coherent condensate.
It is the action carried by the first nontrivial loop sector, up to the
factor $2\pi$.
In this sense $S_0$ is \emph{pre-quantum}: it is an action scale
selected by the theory's own topology, but it has not yet been shown to
be the same quantity that appears in all observed quantum processes.

\paragraph{Why $S_0=\hbar$ is not yet established.}
The structural existence of a distinguished loop action scale does not
yet prove that:
\begin{enumerate}[label=(\roman*)]
  \item all coherent amplitudes are weighted by the same unit $S_0$;
  \item canonical quantisation rules are governed by $S_0$;
  \item the wave-dynamical limit of the coherent branch reproduces the
    observed quantum kinematics with $S_0$ in the r\^ole of $\hbar$;
  \item the numerical value of $S_0$ matches the observed
    $\hbar=1.054\ldots\times10^{-34}\,\mathrm{J\,s}$.
\end{enumerate}
These are additional statements; none follows merely from
the structural existence of nontrivial loop sectors.

\paragraph{What is required for the identification $S_0=\hbar$.}
To identify $S_0$ with the observed Planck constant, the following
multi-step program is required:
\begin{enumerate}[label=(\roman*)]
  \item derive the reduced loop action explicitly within the coherent
    EFT, rather than only introducing its structural
    form~(Appendix~\ref{app:Sloop_calc});
  \item show that the resulting action scale controls the effective
    phase weight of coherent histories in the path-integral sense;
  \item demonstrate that the canonical and wave-dynamical limits of the
    coherent sector are governed by the same scale;
  \item perform the numerical matching to the observed $\hbar$.
\end{enumerate}
This is a multi-step closure problem of the Quantum Sector, not a
one-line identification.

\paragraph{Role of Appendix~\ref{app:Sloop_calc}.}
Appendix~\ref{app:Sloop_calc} explains why the explicit evaluation of
$S_0$ was postponed in ECT basics.
The point is not that the scale is absent, but that its computation
requires a separate reduced loop construction: integration over
transverse loop structure, the effective 1D loop functional, and the
extraction of the minimal topological action contribution.
This is more than reading off a coefficient from the 4D EFT.

\paragraph{Why this is already a major structural result.}
Even before the final identification with $\hbar$, the emergence of
$S_0$ is conceptually significant.
Most formulations of quantum theory begin by postulating a universal
quantum of action.
ECT, by contrast, reaches an action scale through the topology of the
coherent condensate itself.
The theory has therefore already explained \emph{why an action quantum
candidate should exist at all}, even if it has not yet completed the
empirical identification.

\paragraph{Bridge to wave dynamics.}
Once a distinguished action scale exists, the natural next question is
how coherent sectors combine into an effective wave description and
whether the Wick rotation $w\to -ic_*t$ carries the phase-sector
action into a standard kinetic term with $S_0$ as the effective
action unit.
That question leads directly to the wave-dynamical development of the
Quantum Sector, in which the transition from the Euclidean phase action
to Schr\"odinger-like dynamics is made explicit.

\paragraph{Status summary.}
\textit{Established~(Level~A):}
the coherent-topological branch of ECT contains a distinguished
pre-quantum action scale $S_0$ associated with the first nontrivial
loop sector.
\textit{Not yet established~(Level~B, open):}
the full dynamical, canonical, and numerical identification
\begin{equation}
  S_0 = \hbar .
  \label{eq:qs_S0_identification}
\end{equation}
That step belongs to the constructive program of the Quantum Sector.

\paragraph{Summary.}
ECT already contains more than a generic broken-phase field theory:
it contains a topologically selected candidate for a universal action
scale.
This is the correct starting point for the emergence of quantum
kinematics.
The theory remains methodologically honest: $S_0$ is structurally
present before it is empirically identified with $\hbar$.
The next task is to examine whether the coherent phase action develops
into wave dynamics and canonical structure governed by this same scale.

% ============================================================
\subsection{Wave dynamics, Wick rotation, and
  Schr\"odinger-type evolution}
\label{sec:qs_wave}
% ============================================================

\textit{Status: Level~A for the existence of the coherent phase
action~\eqref{eq:qs_phase_action} and for the structural role of the
analytic continuation $w\to-ic_*t$ in the broken phase.
Level~B for the full Schr\"odinger identification, which additionally
requires the explicit identification $S_0=\hbar$ and the final
normalisation of the coherent wave functional.
Connection to ECT basics: Sections~\ref{sec:hbar_status},
\ref{sec:coh_branch}; Appendix~\ref{app:Sloop_calc}.}

\medskip
The previous two subsections established that the coherent-topological
branch is organised by the phase field $\theta$, its winding sectors,
and the distinguished pre-quantum action scale $S_0$.
The question now is whether this coherent sector admits a
wave-dynamical description.
This is the first place in the Quantum Sector where the ECT structure
begins to resemble ordinary quantum theory.
The resemblance must be built, not assumed.

\subsubsection*{From phase action to coherent amplitudes}

The natural starting point is the reduced coherent-phase
action~\eqref{eq:qs_phase_action}, recalled here schematically:
\[
  S_\theta
  = \frac{K_\theta^{\rm eff}}{2}
    \int d^4X\; g_{\rm eff}^{AB}\,\partial_A\theta\,\partial_B\theta
    + \cdots
\]
Because the coherent branch is organised by loop sectors weighted by
their action, the natural amplitude associated with a coherent
history~$\theta(\cdot)$ is the phase weight
\begin{equation}
  \label{eq:qs_weight_S0}
  \mathcal{A}[\theta]
  \;\sim\;
  \exp\!\left(\frac{i\,S_\theta[\theta]}{S_0}\right).
\end{equation}
Here $S_0$ is the distinguished loop action scale~\eqref{eq:qs_Sloop},
and the symbol `$\sim$' indicates that this is a structural assignment,
not yet a normalised path-integral measure.

At this stage Eq.~\eqref{eq:qs_weight_S0} is not the final quantum
path integral.
Its role is more modest and precise: it expresses the fact that once
coherent sectors are classified by their action, the natural language
of their interference is a phase weight normalised by the intrinsic
action scale of the coherent branch.
This is a Level~B statement because the full normalisation requires the
$S_0=\hbar$ identification~\eqref{eq:qs_S0_identification}.

\subsubsection*{Why Wick rotation is necessary and not arbitrary}

The ordered condensate is constructed from a Euclidean arena
$\mathcal{M}^4$, but the broken phase develops a distinguished
Lorentzian direction through the sign structure of the effective kinetic
tensor $g_{\rm eff}^{AB}$.
As established in ECT basics and used throughout Macroscopic Physics,
the broken phase admits the continuation
\begin{equation}
  \label{eq:qs_wick}
  w \;\to\; -\,i\,c_*\,t,
\end{equation}
which transforms the broken-phase kinetic operator into its Lorentzian
hyperbolic form and reconstructs causal propagation.
(The same continuation was used to derive the Lorentzian metric
in Section~\ref{sec:lorentz_order_field} and to develop
cosmological dynamics in Section~\ref{sec:mp_cosmo}.)

In the Quantum Sector the role of Eq.~\eqref{eq:qs_wick} is
complementary: it converts the coherent Euclidean phase action into a
Lorentzian phase-evolution law for time-dependent coherent amplitudes.
This is not an arbitrary formal trick.
It is the same continuation that already connects the Euclidean
ordering arena to observable Lorentzian physics in the geometric branch.

\paragraph{Two complementary uses of the same Wick rotation.}
The emergence of time in ECT therefore has two parallel aspects:
\begin{enumerate}[label=(\roman*)]
  \item in Macroscopic Physics, $w\to -ic_*t$ yields Lorentzian
    geometry, causal cones, and metric-sector propagation;
  \item in the Quantum Sector, the same continuation yields coherent
    temporal phase evolution and therefore wave dynamics.
\end{enumerate}
This parallelism is structurally important: the two branches are not
unrelated, but two uses of the same broken Lorentzian ordering
structure.

\subsubsection*{Lorentzian form of the coherent phase dynamics}

Under the continuation~\eqref{eq:qs_wick}, the phase-sector
action~\eqref{eq:qs_phase_action} acquires the Lorentzian form
\begin{equation}
  \label{eq:qs_phase_action_L}
  S_\theta^{(L)}
  \;\sim\;
  \frac{K_\theta^{\rm eff}}{2}
  \int dt\,d^3x
  \left[
    \frac{1}{c_*^2}(\partial_t\theta)^2
    - (\nabla\theta)^2
  \right]
  + \cdots
\end{equation}
This is again a Level~B effective closure.
Its importance is that it provides the correct Lorentzian kinetic
structure for coherent phase evolution, with the spatial gradient and
time-derivative terms separated by the effective speed~$c_*$.

\subsubsection*{Schr\"odinger-type reduction}

\paragraph{The WKB-like ansatz.}
To pass from the second-order Lorentzian phase dynamics of
Eq.~\eqref{eq:qs_phase_action_L} to an effective first-order
evolution law for a complex coherent amplitude, one introduces the
coherent wave functional (a WKB-like polar decomposition of the
order parameter; cf.~Eq.~\eqref{eq:qs_polar}):
\begin{equation}
  \label{eq:qs_wave_ansatz}
  \Psi(x,t)
  =
  \sqrt{\rho(x,t)}\;
  \exp\!\left(\frac{i\,\mathcal{S}(x,t)}{S_0}\right),
\end{equation}
where $\mathcal{S}$ is the effective coherent phase along the history
and $\rho$ is the coarse-grained coherent density.
This ansatz is Level~B: it assumes slow variation of $\rho$ and
$\mathcal{S}$ compared to the microscopic scale.

\paragraph{Emergence of Schr\"odinger-type equation.}
Inserting~\eqref{eq:qs_wave_ansatz} into the reduced coherent dynamics
and retaining the leading slowly-varying terms yields, schematically,
a first-order coherent evolution law:
\begin{equation}
  \label{eq:qs_schrodinger_type}
  i\,S_0\,\partial_t\Psi
  =
  \hat{H}_{\rm coh}\,\Psi
  + \cdots
\end{equation}
where $\hat{H}_{\rm coh}$ is the effective coherent Hamiltonian
(arising from the spatial gradient and potential terms in the reduced
phase action), and the ellipsis denotes decoherence, environmental,
and higher-gradient corrections.

The appearance of $S_0$ in the r\^ole of $\hbar$ in
Eq.~\eqref{eq:qs_schrodinger_type} is structurally natural: $S_0$ is
the distinguished action unit of the coherent branch, and it is
$S_0$---not an externally postulated constant---that sets the scale
of the effective coherent dynamics.

\paragraph{Status of Eq.~\eqref{eq:qs_schrodinger_type}.}
This is \emph{not} yet the final observed Schr\"odinger equation.
What is structurally established is the following: once the coherent
branch possesses (i)~a distinguished action scale $S_0$,
(ii)~a Lorentzian continuation~\eqref{eq:qs_wick}, and (iii)~a reduced
phase action for smooth coherent configurations, then a
Schr\"odinger-type first-order coherent wave description becomes the
natural effective language of the low-energy coherent regime.

What is still missing for the full identification with ordinary
quantum mechanics:
\begin{enumerate}[label=(\roman*)]
  \item the explicit derivation of $\hat{H}_{\rm coh}$ for the
    relevant coherent sector;
  \item the identification $S_0=\hbar$~\eqref{eq:qs_S0_identification};
  \item the control of decoherence and environment
    terms~(Section~\ref{sec:qs_decoherence}).
\end{enumerate}
For these reasons Eq.~\eqref{eq:qs_schrodinger_type} is currently a
Level~B emergence statement.

\subsubsection*{Why this is already structurally nontrivial}

Even before the final $S_0=\hbar$ identification, ECT has gone beyond a
generic broken-phase field theory.
The theory organises the coherent branch by:
\begin{itemize}
  \item topological loop sectors with integer winding~(P5);
  \item an intrinsic action scale $S_0$ determined by the topology;
  \item Lorentzian time reconstruction through the same broken-phase
    continuation used in Macroscopic Physics;
  \item a natural first-order coherent wave reduction governed by $S_0$.
\end{itemize}
This is precisely the structural content required from a theory that
aims to explain \emph{why wave mechanics exists at all}, rather than
postulating it from the outset.

\subsubsection*{Bridge to canonical structure}

The logical sequence of the Quantum Sector is:
\begin{equation}
  \label{eq:qs_logic_chain}
  \text{phase sectors}
  \;\longrightarrow\;
  S_0
  \;\longrightarrow\;
  \text{wave dynamics}
  \;\longrightarrow\;
  \text{canonical structure}.
\end{equation}
Once a Schr\"odinger-type coherent evolution law is available, the next
question is whether the coherent variables admit a canonical
organisation and whether their brackets or commutators are normalised
by the same scale $S_0$.
That question belongs to Section~\ref{sec:qs_ccr}.

\paragraph{Status summary.}
\textit{Established~(Level~A):}
the coherent branch admits a Lorentzian kinetic structure after the
continuation~\eqref{eq:qs_wick}, and this structure leads naturally
to a first-order coherent wave reduction with $S_0$ as the action unit.
\textit{Not yet established~(Level~B):}
the full identification of that reduction with ordinary quantum
mechanics in normalised form~(requires $S_0=\hbar$ and
$\hat{H}_{\rm coh}$ derived explicitly).

\paragraph{Summary.}
The coherent-topological branch of ECT does not jump from loop
quantisation to operator axioms.
The logic is controlled: the reduced phase action, together with the
broken-phase Wick continuation~\eqref{eq:qs_wick} and the distinguished
action scale $S_0$, naturally produces a Schr\"odinger-type wave
description of smooth coherent configurations.
This prepares the next step: the canonical structure of the coherent
sector and the status of commutation relations.
